\bookchapter{פרק 6: הגן על חייך הדיגיטליים – אבטחה ופרטיות}

\chapterintro{אבטחת טלפון = אבטחת פרטיות + אבטחת רכוש. בעידן שבו נתונים הם נכס, להגן על הטלפון זה להגן על עצמך.}

\booksection{נעילת מסך: קו ההגנה הפיזי הראשון}

סיסמה / קוד PIN

\booksubsection{דירוג רמת אבטחה:}

\begin{enumerate}
  \item \textbf{זיהוי טביעת אצבע}: איזון מושלם בין אבטחה לנוחות.
  \item \textbf{זיהוי פנים}: הנוח ביותר, אך רמת האבטחה תלויה בטכנולוגיה (אור מובנה 3D בטוח, תמונה 2D קלה לזיוף).
  \item \textbf{ביטול נעילת דפוס}:נוח אך קל להשאיר עקבות שניתן לנחש.
  \item \textbf{הגדרת נעילת מסך חזקה:}הקפד להגדיר סיסמה מורכבת (אל תשתמש ביום הולדת, 123456).
\end{enumerate}

\booksubsection{הגדרת נעילת מסך חזקה：}

\begin{itemize}
  \item הגדר את זמן נעילת המסך ל"מיידי" או הקצר ביותר.
  \item 2. "מצא את הטלפון שלי": כלי אולטימטיבי נגד גניבה
  \item חובה להפעיל!
\end{itemize}

\booksection{“מצא את הטלפון שלי”: כלי אולטימטיבי נגד גניבה}

\textbf{איתור}:איפה שהטלפון אבד, חפש במפה.

\begin{itemize}
  \item \textbf{השמעת צליל}:לגרום לטלפון לצלצל, מתאים לחיפוש בבית.
  \item \textbf{מצב אובדן}:נעל את הטלפון מרחוק והצג מספר ליצירת קשר.
  \item \textbf{מחיקה מרחוק}:המוצא האחרון, נקה את כל הנתונים (אך לאחר המחיקה לא ניתן לאתר יותר).
  \item \textbf{עבור אל [הגדרות], באייפון חפש "מצא את ה-iPhone שלי", באנדרואיד חפש "מצא את המכשיר שלי" או שירות דומה של המותג, וודא שהוא מופעל (ירוק).}3. עדכוני מערכת: התקן תיקונים, סתום פרצות
\end{itemize}

המטרה העיקרית של עדכוני מערכת היא לתקן פרצות אבטחה, בדיוק כמו לתקן חור בקיר הבית. לכן, כשרואים התראת עדכון מערכת, כדאי לשדרג בהקדם האפשרי.

\booksection{עדכוני מערכת: התקן תיקונים, סתום פרצות}

4. זיהוי ומניעת הונאות טלפוניות

בעידן של התפתחות מהירה בטכנולוגיית AI, תפיסת “לראות זה להאמין” המסורתית כבר אינה תקפה. נוכלים מנצלים טכנולוגיות AI כמו Deepfake, והפכו מהונאות “רשת רחבה” לפשעי “ציד ממוקד” משודרגים.

\booksection{זיהוי ומניעת הונאות טלפוניות}

(1) הונאת החלפת פנים וחיקוי קול באמצעות AI

\booksubsection{סוגי ההונאות הנפוצים ביותר בעידן ה-AI}

\textbf{תאונת דרכים פתאומית, הלוואה דחופה, חטיפה}\par

וכו'. הונאות אלו מנצלות קשרים משפחתיים ורגשיים,\textbf{והן מטעות ביותר}.\textbf{(2) הונאת התחזות לרשות / שירות לקוחות}: טכנולוגיית ה-AI מסוגלת לחקות במדויק את הקול והטון של שירות לקוחות רשמי ושל אנשי אכיפה, ואף לייצר צילומי מסך כוזבים של העברות, מסמכים רשמיים אלקטרוניים וצווי מעצר ריאליסטיים במיוחד. באמצעות\textbf{יצירת פחד או לחץ סמכותי}היא מפתה את הקורבן להעביר כספים או להפעיל שיתוף מסך, וכך גונבת את פרטי החשבון.

\textbf{יצירת פחד או לחץ מסמכות}\par

כדי לגרום לקורבן להעביר כסף או להפעיל שיתוף מסך, וכך לגנוב פרטי חשבון.\textbf{(3) הונאת השקעות באמצעות AI}שימוש באלגוריתמי AI לזיוף צילומי מסך של רווחים “בטוחים”, ממשקי תוכנות בחירת מניות וירטואליות או דוחות ניתוח השקעות “ברמת מומחה”, תוך שימוש בהבטחות לתשואות גבוהות כפיתיון למכירת שירותי השקעה יקרים או לגרום לקורבן להעביר כסף לחשבונות מסחר שבשליטתם,

\textbf{ובסוף לברוח עם הכסף}\par

\textbf{2. דרך לשבור את המצור – גישת ההגנה של “מראת הקסם”}: ה-AI יכולה לשכפל קול ופנים, אך גם היא אינה חפה מפגמים — יש כמה נקודות תורפה שאפשר לנצל:

\booksubsection{דרך לשבור את המצור – גישת ההגנה של “מראת הקסם”}

:במהלך שיחת וידאו, בקש מהצד השני לבצע סדרת תנועות מוגדרות מראש, כמו “העבר את היד במהירות מול הפנים”, “הצג את קווי המתאר של האוזן בפרופיל” או “מצמץ 5 פעמים ברציפות ובמהירות”. טכנולוגיית החלפת הפנים הנוכחית ב-AI, כאשר מתמודדת עם חסימות דינמיות מורכבות, תנועות מהירות ושינויי תאורה,

\begin{itemize}
  \item \textbf{נוטה מאוד לגמגום תמונה, עיוות פנים או חדירת קצוות}.\textbf{שיטת אימות בסיסמה סודית}:קבע והשתמש בשאלות אימות סודיות מוגדרות מראש, כמו "מה הזמנת בארוחה המשותפת האחרונה שלנו?" או "מה הייתה המתנה הראשונה שנתתי לך?" – פרטים שרק שני הצדדים יודעים. אם הצד השני מהסס או עונה תשובה לא נכונה, יש
  \item \textbf{לנתק מיד ולהפעיל תוכנית חירום}.
  \item \textbf{שיטת האימות הכפול}: הקפידו על "ארבעת הלאווים" — אל תאמינו בקלות, אל תחשפו, אל תעבירו כספים, אל תעבירו מידע רגיש; בעסקאות כספיות\textbf{חייגו בחזרה למספר המקורי המוכר או אמתו פנים אל פנים}.
  \item \textbf{ארבעת הלא ואימות אחד}" שפרסם המשרד לביטחון הציבור – אל תאמין להודעות מאנשים זרים, אל תעביר כספים בפזיזות, אל תדליף מידע אישי רגיש, אל תפיץ תוכן חשוד. כאשר מדובר בעסקאות כספיות, יש\textbf{לחייג בחזרה למספר המקורי הידוע או לאשר פנים אל פנים}. בעת ביצוע העברה בנקאית, מומלץ לבחור במצב "העברה תוך 24 שעות" כדי להשאיר חלון זמן מספק לזיהוי הונאה אפשרית ולהצלת הכספים.\textbf{3. תרחישים מסורתיים נפוצים (עדיין יש להיזהר):}זכייה בפרס / החזרים על ביקורות מזויפות
\end{itemize}

\booksubsection{תרחישים מסורתיים נפוצים (עדיין יש להיזהר)：}

\begin{itemize}
  \item \textbf{הונאת רומנטיקה מקוונת}:"שחיטת חזיר", בונים קשר רגשי ואז מרמים כסף.
  \item \textbf{התחזות למכר}:בקשת הלוואה לאחר גניבת חשבון (קשה עוד יותר לזהות בעידן ה-AI).
  \item \textbf{4. הגנה מובנית בטלפון:}חסימת שיחות מטרידות.
\end{itemize}

\booksubsection{הגנה מובנית בטלפון：}

\begin{itemize}
  \item בדיקת אבטחת קישורים בהודעות SMS.
  \item 5. נוסחת מניעת הונאות:
  \item אל תאמין בקלות, אל תחשוף, אל תלחץ, אל תעביר כסף, אם יש ספק – ברר קודם!
\end{itemize}

\booksubsection{נוסחת מניעת הונאות：}

\textbf{5. הגנה מפני “Doxing” והתקפות הנדסה חברתית: ניתוק שרשרת דליפת הפרטיות}העיקרון המרכזי של “Doxing” (התקפת הנדסה חברתית) הוא

\booksection{הגנה מפני “Doxing” והתקפות הנדסה חברתית: ניתוק שרשרת דליפת הפרטיות}

. המפתח להגנה הוא קביעת קווי הגנה ברורים לפרטיות האישית ובניית מערכת הגנה רב-שכבתית.\textbf{1. טיהור תמונות ממידע רגיש}לפני שיתוף תמונות ברשתות חברתיות, יש להשתמש בכלי עריכת תמונות מקצועיים

\booksubsection{טיהור תמונות ממידע רגיש}

שיטות בסיסיות לניקוי EXIF נמצאות בפרק 2 ״ניהול תמונות ומערכת קבצים״. לפני שיתוף תמונות ברשתות חברתיות, יש צורך בהמשך\textbf{מספרי בתים, לוחיות רישוי, ברקודים של כרטיסי נסיעה, מספרי תעודת זהות}וכדומה, כדי להימנע מחשיפה לא מכוונת של מסלול תנועה אישי.\textbf{iOS:[הגדרות] > [פרטיות ואבטחה] > [מצלמה] ניתן לכבות “מיקום מדויק”; בעת שיתוף, בחירה ב“תמונה נוכחית בלבד” מסירה מטא-נתונים.}אנדרואיד:חלק מאפליקציות הגלריה (כמו Google Photos) מציגות התראה בעת שיתוף האם לשמור מידע מיקום.

\begin{itemize}
  \item כלים אוניברסליים:ExifTool, Metapho וכו' מאפשרים צפייה/עריכה של EXIF.
  \item 2. צמצום הרשאות אפליקציות
  \item בדוק באופן קבוע וסגור הרשאות לא הכרחיות כמו
\end{itemize}

\booksubsection{צמצום הרשאות אפליקציות}

שיטות ניהול הרשאות בסיסיות בפרק 4 ״ניהול הרשאות״. בסיטואציית ״מניעת חשיפת מידע״, יש צורך נוסף לבדוק ולסגור באופן תקופתי הרשאות של אפליקציות מסחר וכלים\textbf{כשרואים בקשה ל“הרשאת אנשי קשר”, שאל קודם: האם האפליקציה הזו באמת צריכה לדעת את כל אנשי הקשר שלי?}אפליקציות שלא השתמשת בהן שלושה חודשים – הסר מיד. הרשאות צריכות להצטמצם עם הזמן, לא להתרחב.

\begin{itemize}
  \item 3. אבטחת מכשירים חכמים
  \item מצלמות חכמות, מנעולים חכמים, רמקולים חכמים ומכשירי IoT אחרים – בעת השימוש הראשון יש
\end{itemize}

\booksubsection{אבטחת מכשירים חכמים}

ולהגדיר סיסמה חזקה וייחודית הכוללת אותיות גדולות וקטנות, מספרים ותווים מיוחדים; יש להשאיר פונקציות גישה מרחוק שאינן בשימוש תכוף כבויות, לבדוק באופן קבוע עדכוני קושחה ולתקן פרצות אבטחה בזמן,\textbf{כדי למנוע מפורצים לחדור ולצפות במרחב הפרטי של הבית}.

\begin{tcbwarn}{אזהרת מקרה אמיתי}
ב-2022 דיווחה משטרה מקומית: האקרים פיצחו את סיסמת ברירת המחדל של מצלמות ביתיות ממותגות (משתמשים רבים מעולם לא שינו אותה), גנבו צילומי חדרי שינה של אלפי משפחות ומכרו אותם ברשת.\textbf{אי-שינוי סיסמת ברירת המחדל = להשאיר את דלת הבית פתוחה לרווחה}.

\end{tcbwarn}

\booksection{הגנת פרטיות: הנתונים שלך הם הרכוש שלך}

התמונות, יומני הצ'אט ומסלולי המיקום בטלפון שלך הם בבחינת “הנכסים הדיגיטליים” שלך. ברגע שהם דולפים — במקרה הקל תגיע הטרדת פרסומות, ובמקרה הכבד חשיפה, הונאה ואף סחיטה.
